\section*{Апробация системы}

Апробация разработанной системы для анализа мультиагентных алгоритмов определения позиций объектов проводилась в рамках научно-исследовательской деятельности и включала представление результатов на научных мероприятиях, а также их публикацию в рецензируемых изданиях.

Основные результаты работы были представлены в виде доклада на XXVIII конференции молодых ученых «Навигация и управление движением» \cite{youngSc}, организованной АО «Концерн ЦНИИ \textit{Электроприбор}» \cite{elektropribor}. Доклад был сделан в рамках секции «Прикладные задачи навигации и управления движением» и охватывал вопросы разработки архитектуры системы, реализации симуляционной среды, а также интеграции и сравнительного анализа мультиагентных алгоритмов. В ходе выступления были продемонстрированы возможности системы по проведению воспроизводимых вычислительных экспериментов, гибкой настройке параметров моделирования и визуализации результатов.


Кроме того, по теме исследования был подготовлен реферат, который прошёл рецензирование и был принят к публикации в материалах указанной конференции. В публикации отражены ключевые положения работы, включая постановку задачи, описание архитектуры системы и результаты экспериментальных исследований.

Работа также была представлены на конференции «СПИСОК‑2026» \cite{spisok} в секции «Комплексная автоматизация БПЛА: от планирования полётного задания до анализа данных».


Разработанная система была использована в качестве инструмента для проведения серии вычислительных экспериментов, направленных на сравнительный анализ алгоритма на основе SPSA с консенсусом и его ускоренной версии. Проведённые эксперименты включали пакетный запуск алгоритмов с помехами. Это позволило оценить сходимость алгоритмов.

Полученные результаты показали, что ускоренная версия алгоритма обладает более высокой скоростью сходимости при сохранении устойчивости и точности оценивания, что подтверждает эффективность предложенных модификаций. Система обеспечила удобную платформу для проведения анализа за счёт поддержки пакетного запуска экспериментов и автоматизированного сравнения результатов.

Результаты проведённого исследования, полученные с использованием разработанной системы, были опубликованы в статье в журнале \textit{IEEE Transactions on Automatic Control} \cite{Accelerated_SPSA}, что подтверждает их научную значимость и новизну. 

На текущий момент по работе была написана полноценная статья, описывающая систему, и подана на ApPLIED Workshop \cite{Applied} (в рамках PODC-2026, Эгхем, Великобритания, 6 июля 2026 г.).

Таким образом, разработанная система прошла апробацию как в рамках научных конференций, так и в процессе подготовки публикаций в высокорейтинговых научных изданиях.